\subsection{The original realization: closed programs and saturated types}
\label{sec:p07-ch88-realization}

This is the realization argument of the original article's sections~4--5,
in the corrected binding/equality presentations already declared above.
It is not the annotation-retaining interpretation used for full annotated
normalization. Programs, semantic constructor values and realized valuations
are different objects \cite[pp.~105--110]{CoquandHuet1988}.

\begin{hbdefinition}{def:p07-ch88-native-saturation}{The source closed-program saturated sets}
Let \(\mathcal Y\) be the set of finite closed untyped lambda terms with
an inert constant \(\Omega\). Here \(\Omega\) is a symbol with no
reduction rule, not the usual divergent self-application abbreviation.
Use ordinary compatible beta reduction. A saturated set \(A\subseteq
\mathcal Y\) contains \(\Omega\), contains every
\(\Omega b_1\cdots b_n\) with strongly normalizing closed arguments,
and consists only of strongly normalizing terms. It has the source's
head-expansion closure: if \(b\) is strongly normalizing and
\((u[b/x])b_1\cdots b_n\in A\), then
\(((\lambda x.u)b)b_1\cdots b_n\in A\), whenever these whole
expressions are closed. The template \(u\) may have the one free variable
\(x\); the argument and spine terms are closed. Write \(\mathcal U\)
for all such saturated sets. No reduction-closure or general
beta-equivalence-invariance of membership is added.

For \(A\in\mathcal U\) and a total family \(F:\mathcal Y\to
\mathcal U\), define
\[
 \Pi(A,F)=\{t\in\mathcal Y:\forall b\in A,\ tb\in F(b)\}.
\]
This dependent product uses the same actual argument \(b\) in the
application and the family index. It is not an ordinary full function
space of set-theoretic functions \cite[p.~107]{CoquandHuet1988}.
\end{hbdefinition}

\begin{hblemma}{thm:p07-ch88-native-saturation}{proved=annex:p7.ch88.native-saturation}{The source closure and introduction properties}
\(\mathcal U\) is nonempty and is closed under nonempty intersections
and the displayed dependent product. If \(u[b/x]\in F(b)\) for
every \(b\in A\), with the same one-variable template \(u\), then
\(\lambda x.u\in\Pi(A,F)\). These properties do not assume source
term normalization.
\end{hblemma}

\begin{hbdefinition}{def:p07-ch88-native-functionality}{Functionality and the two ambient models}
For variant \(q\in\{r,f\}\), an object \(M\in\mathcal A_i\)
has functionality \(\varphi_q(M)=\mathcal Y\), and
\(\varphi_q(\ast)=\mathcal U\). For a logical context
\([x:A]\Delta\), the restricted model \(r\) uses all functions
\(\varphi_r(A)\to\varphi_r(\Delta)\). In the full model \(f\),
use all such functions when \(A\) is a logical context; when \(A\)
is an object, use precisely the functions
\(g:\mathcal Y\to\varphi_f(\Delta)\) satisfying
\(b=_\beta c\Rightarrow g(b)=g(c)\).
The index \(q\) explicitly qualifies the source's two choices.

For an environment \(\Gamma=[x_n:A_n]\cdots[x_1:A_1]\ast\),
the source also uses functionality for the ambient Cartesian product.
To distinguish this role from a logical-context term's function space,
write
\(\varphi_{q,\mathrm{env}}(\Gamma)=
\varphi_q(A_n)\times\cdots\times\varphi_q(A_1)\).
An ambient valuation is a tuple \(e=(a_n,\ldots,a_1)\) in that
product. A coordinate whose annotation is a logical context is a
semantic constructor value; one with an object-type annotation is a
closed program. These ambient values need not yet realize their types.
All domains are ordinary sets built by finite products/function spaces
from \(\mathcal Y\) and \(\mathcal U\). They are nonempty:
\(\Omega\) inhabits \(\mathcal Y\), the set of closed SN terms
inhabits \(\mathcal U\), and constant functions give the recursive
inhabitants, including the full beta-invariant function spaces
\cite[pp.~107--108]{CoquandHuet1988}.
\end{hbdefinition}

\begin{hbdefinition}{def:p07-ch88-native-values}{The actual two-level realization}
Write \(\rho_{q,\Gamma}(M)(e)\) for the source interpretation.
If \(M\)'s type is an object, take the stripped program
\(\nu_\Gamma(M)\) and substitute the program coordinates selected
by \(j_\Gamma\). Thus the program value is in \(\mathcal Y\) and
is independent of constructor coordinates. If \(M\)'s type is a
logical context, its constructor interpretation is defined by the following
structural clauses. A variable selects its coordinate. A constructor
lambda curries the body value over the coordinate of its domain.
A constructor application applies its function value to the actual
argument value. For a small type \([x:A]B:\ast\), put
\[
\rho_{q,\Gamma}([x:A]B)(e)=
\begin{cases}
\Pi(\rho_{q,\Gamma}(A)(e),
 b\mapsto\rho_{q,\Gamma[x:A]}(B)(e,b)),&A\text{ an object},\\
\displaystyle\bigcap_{a\in\varphi_q(A)}
 \rho_{q,\Gamma[x:A]}(B)(e,a),&A\text{ a logical context}.
\end{cases}
\]
The first line quantifies over the realized object domain; the second
over every semantic constructor value. They are different operations.

Interpretations are defined on well-sorted, role-correct endpoints.
The source-image structural typing certificates provided by
\cref{thm:p07-ch88-full-transport} certify the constructor/proof modes,
annotation scopes and function arities; the accompanying structural CC
facts use no normalization. They are not an inverse adequacy assertion.
In particular typed equality alone need not supply a typing of an
arbitrary product endpoint. Conversion uses the sorted endpoints established
by that transport. The clauses do not define an interpretation of a
logical-context term as an object inhabitant
\cite[\S5.2, pp.~108--109]{CoquandHuet1988}.
\end{hbdefinition}

\begin{hblemma}{thm:p07-ch88-native-values}{proved=annex:p7.ch88.native-values}{Defined functionalities, scoping and equality compatibility}
The constructor/proof values just defined are total on their ambient
valuation domains and belong to the functionality of their prescribed
type. Small-type values belong to \(\mathcal U\).
Relocation selects the same older coordinates and interpretation commutes
with a typed binder-removal substitution by inserting the actual value
of its argument in the valuation.

For interpretable typed-equality endpoints, functionality spaces agree.
Restricted constructor equality gives literally equal values. In the full
variant, constructor/type equality also gives literally equal values,
whereas proof equality gives beta-equal stripped programs. This last
conclusion does not assert equal saturated-set membership under arbitrary
beta equivalence. The full model's beta-invariant function requirement
is used when a constructor receives beta-equal proof arguments.
\end{hblemma}

\begin{hbdefinition}{def:p07-ch88-native-realizers}{The source realized subenvironment}
Define \(D_q(\ast)=\{()\}\). For an object-type annotation \(A\),
\[
 D_q(\Gamma[x:A])=
 \{(e,b):e\in D_q(\Gamma),\ b\in\rho_{q,\Gamma}(A)(e)\}.
\]
For a logical-context annotation \(C\), put
\(D_q(\Gamma[x:C])=D_q(\Gamma)\times\varphi_q(C)\).
Thus \(D_q(\Gamma)\subseteq\varphi_{q,\mathrm{env}}(\Gamma)\).
The dependency of an object coordinate on its older prefix is retained;
this is not the whole ambient Cartesian product. For example
\(D([A:\ast][x:A]\ast)=\{(R,b):R\in\mathcal U, b\in R\}\)
\cite[\S5.3, p.~109]{CoquandHuet1988}.
\end{hbdefinition}

\begin{hbtheorem}{thm:p07-ch88-native-soundness}{proved=annex:p7.ch88.native-soundness}{The complete original proof-level realization argument}
For either corrected source variant with its corresponding model, if
\(\Gamma\vdash M:P\) and \(P\in\mathcal A_i^{|\Gamma|}\),
then for every \(e\in D_q(\Gamma)\),
\[
 \rho_{q,\Gamma}(M)(e)\in\rho_{q,\Gamma}(P)(e).
\]
The right side is a saturated set by the source-image sort certificate,
and the left a closed pure program. In particular this proves the original
two-premise statement when \(\Gamma\vdash P:\ast\) is also given.
This expands the source's stated structural induction; it uses neither
the full annotated normalization conclusion nor the annotation-retaining
interpretation \cite[\S5.4, p.~109]{CoquandHuet1988}.
\end{hbtheorem}

\begin{hbcorollary}{thm:p07-ch88-native-closed-image}{proved=annex:p7.ch88.native-closed-image}{The actual closed stripped-image normalization}
If \(\vdash M:P\) and \(P\in\mathcal A_i^0\), then
\(\nu_\ast(M)=\rho_{q,\ast}(M)(())\) is strongly normalizing.
This is the source's closed program-image conclusion, not normalization
of every retained annotation of \(M\) \cite[p.~110]{CoquandHuet1988}.
\end{hbcorollary}

\begin{hbtheorem}{thm:p07-ch88-native-consistency}{proved=annex:p7.ch88.native-consistency}{The source's uninhabited proposition}
The small type \([A:\ast]A\) has no closed source proof in either
corrected variant. The inert marker belongs to the semantic interpretation
but cannot be created by beta reduction of a source-stripped program
\cite[Corollary~2, p.~110]{CoquandHuet1988}.
\end{hbtheorem}

\begin{hbmethod}{def:p07-ch88-source-checking}{The actual source checking extent}
The p.~104 relation has the logical typing guards and explicit transitivity
already displayed on its sheet. On a constructed logical redex the source
assigns complexity \(\delta(P)\) to its domain and collects the
multiset of logical-redex complexities. It asserts a decrease for an
innermost logical contraction, then describes type synthesis as an
induction on the sum of term and environment lengths. These are source
proof outlines, not a complete algorithm specification in this subsection.
They require the logical-redex classification, rank stability, preservation
of the actual typed regime, and the smaller-rank newly created redex argument.
Innermost here is not an instruction to count every proof-level beta redex.
The full calculus's checking claim additionally refers to normalization of
all constructed terms. The p.~112 normal-type presentation normalizes
declaration annotations and application result types, after those
metatheorems are available \cite[pp.~104--105,112]{CoquandHuet1988}.
\end{hbmethod}

\begin{proofboundary}{def:p07-ch88-source-checking-boundary}{What realization does not establish about checking}
The realization theorem proves membership and erased-program termination;
it is not a type-inference algorithm or an inverse to the forward CC bridge.
The source does not enumerate the checking procedure's acceptance/failure
branches or its full soundness/completeness proof. A complete original
checker still needs source generation and type-preservation/adequacy for
its actual conversion regime and for the normal-type variant. Those
obligations are not replaced by a successful target typing or by normalizing
two target images. The original rank/innermost discussion and cited
normalization dependencies keep their distinct proof standing.
\end{proofboundary}
